\subsection{Neuron Model and Simulation Environment}
All computational experiments were conducted using the NEST simulator, a framework specifically optimized for the simulation of large-scale spiking neural networks. The fundamental processing units of the network are modeled as leaky integrate-and-fire (LIF) neurons with exact integration and exponential-decaying postsynaptic currents. The subthreshold dynamics of the membrane potential ($V_m$) are governed by the membrane time constant ($\tau_m = 10$ ms), membrane capacity ($C_m = 250$ pF), and a resting potential ($V_{rest}$) of -65 mV. An action potential is generated when $V_m$ reaches a critical threshold ($V_{th} = -50$ mV), after which the potential is strictly clamped to a reset value ($V_{reset} = -65$ mV) for the duration of an absolute refractory period ($\tau_{ref} = 2$ ms). Synaptic parameters were carefully tuned to ensure network stability and to sustain the asynchronous-irregular firing regimes characteristic of awake cortical states. Specifically, excitatory synaptic weights were set to $w = 87.8 \pm 8.8$ pA with a decay time constant $\tau_{syn} = 0.5$ ms, and the relative inhibitory synaptic strength ($g$) was fixed at -4. Transmission delays were uniformly distributed at $1.5 \pm 0.75$ ms for excitatory synapses and $0.75 \pm 0.38$ ms for inhibitory synapses. To emulate background cortical activity, each neuron received external inputs modeled as Poisson spike trains with a mean rate of 8 Hz.

\subsection{Laminar Microcircuit Architecture}
The primary visual cortex (V1) module is designed to represent a $1 \text{ mm}^2$ patch of cortical tissue, based on the full-scale canonical microcircuit model proposed by Potjans and Diesmann (2014). It comprises approximately 77,000 neurons distributed across four distinct laminar layers: L2/3, L4, L5, and L6. Within each layer, the network maintains a physiological balance of specific excitatory and inhibitory populations, approximating an 80/20 ratio. Local recurrent connectivity is established using empirically derived connection probabilities. 

To capture intra-areal spatial integration, long-range horizontal projections within V1 are incorporated. Crucially, these lateral connections are restricted to the superficial L2/3 layer and are functionally segregated based on distance. Nearby connections link functionally similar adjacent columns (iso-orientation) with both excitatory and inhibitory synapses (transmission delay $d_{n.lat} = 5 \pm 1.5$ ms). In contrast, distant connections link columns with disparate orientation preferences and operate via a predominantly suppressive mechanism (disynaptic inhibition) with a slower conduction delay ($d_{d.lat} = 8 \pm 2.5$ ms).

\subsection{Hierarchical Integration and Competitive Dynamics}
To investigate the mechanistic origins of extra-classical receptive field (eCRF) effects, the baseline architecture was extended to include a higher-order visual area (V2). The V2 module was adapted to mirror the specific cytoarchitectural variations in laminar thickness and neuronal densities of the macaque V2 region. Inter-areal communication is mediated by targeted topographic projections operating with fast transmission delays ($2 \pm 0.5$ ms). Feedforward (FF) pathways originate from the superficial layers (L2/3) of V1 and terminate predominantly in L4 and L6 of V2. Conversely, top-down feedback (FB) projections originate from the deep layers (L5 and L6) of V2 and target the superficial (L2/3) and infragranular (L5) layers of V1, anatomically avoiding layer 4 to provide non-linear modulation rather than driving inputs.

Crucially, to test the hypothesis that eCRF phenomena emerge from inter-areal interactions rather than purely local filtering, a lateral competition mechanism was embedded within the V2 module. This architecture utilizes two spatially distinct, feature-aligned V2 modules that engage in lateral competition mediated by strong, reciprocal inhibitory connections. By enforcing this hierarchical winner-take-all competition, the top-down feedback signals are non-linearly refined before modulating V1 activity. This architectural feature is designed to overcome the limitations of simple linear feedback models, ultimately reproducing the robust, stimulus-specific surround suppression observed in physiological recordings.

\subsection{Sensory Stimulation Protocols}
Bottom-up sensory stimulation was emulated through thalamocortical afferents originating from the lateral geniculate nucleus (LGN), projecting primarily to layers L4 and L6 of V1. Thalamic inputs were modeled as Poisson spike trains. To simulate the classical receptive field (CRF) response, the thalamic firing rate directed at the central microcolumn was increased to 20 Hz (simulating optimal stimulus contrast) after a 500 ms baseline period. To evaluate the extra-classical modulations, a dual asynchronous stimulation paradigm was employed: following the central CRF activation at 500 ms, a distant peripheral column was stimulated at 1000 ms. The intensity of this peripheral contextual stimulus was systematically varied (10, 20, and 30 Hz) to quantify the magnitude of surround suppression.

\subsection{Data Analysis and Local Field Potentials}
Network dynamics and mean firing rates were evaluated using peri-stimulus time histograms (PSTH). To connect the spiking network model to macroscopic electrophysiological biomarkers—such as gamma-band oscillations and excitation/inhibition (E/I) balance alterations—we simulated the Local Field Potential (LFP) using a spatiotemporal kernel-based approximation. In this framework, the macroscopic LFP was reconstructed as the linear superposition of unitary field potentials (uLFP) generated by individual action potentials, enabling detailed spectral analysis (PSD) of the network's rhythmic activity under distinct contextual modulations.